\documentclass[11pt]{amsart}
\usepackage{amsmath,amssymb,amsthm}
\usepackage[margin=1.15in]{geometry}
\usepackage[hidelinks]{hyperref}
\newtheorem{theorem}{Theorem}[section]
\newtheorem{corollary}[theorem]{Corollary}
\newtheorem{lemma}[theorem]{Lemma}
\theoremstyle{remark}
\newtheorem{example}[theorem]{Example}
\newcommand{\doi}[1]{\href{https://doi.org/#1}{doi:\nolinkurl{#1}}}
\newcommand{\arxiv}[1]{\href{https://arxiv.org/abs/#1}{arXiv:\nolinkurl{#1}}}
\newcommand{\R}{\mathbb{R}}
\newcommand{\Z}{\mathbb{Z}}
\newcommand{\conv}{\operatorname{conv}}
\newcommand{\spn}{\operatorname{span}}
\newcommand{\supp}{\operatorname{supp}}
\newcommand{\relint}{\operatorname{relint}}
\setlength{\emergencystretch}{2.5em}

\begin{document}

\title[Topology of the recurrent and transient sets]{The topology of the recurrent and transient sets of homogeneous
random walks}
\author{Yitzchak Shmalo}
\address{Einstein Institute of Mathematics, The Hebrew University of Jerusalem, Jerusalem, Israel}
\email{yspennstate@gmail.com}
\subjclass[2020]{60G50, 60J10, 52B11, 55U10}
\keywords{random walk, recurrence, transience, abelian group, path-connected, simply connected, convexity,
nerve theorem}

\begin{abstract}
Each point $a$ of a simplex $A$ defines the random walk on a finitely generated abelian group that takes the $i$-th
step of a fixed finite set $S$ with probability $a_i$. Li, Mossel and Weiss asked when the set $R$ of recurrent walks
is path-connected or simply connected, and when the set $R^c$ of transient walks is path-connected, simply connected or
convex. We answer these questions in terms of the projections $v_i$ of the steps to the free part $\Z^n$ of the group.
Apart from two simple cases, in which $R$ is contractible or convex, $R$ is homotopy equivalent to a bipartite graph
whose vertices are the lines and planes through the origin on which the steps balance. The set $R^c$ is empty when
every $v_i$ is zero, and otherwise contractible unless $0$ lies in the relative interior of the convex hull of the $v_i$ and these vectors span a line or a plane; then
it has two components or is homotopy equivalent to a circle. It is convex exactly when every recurrent walk is
positive recurrent.
\end{abstract}

\maketitle

\section{Introduction}\label{sec:intro}

Let $G$ be a finitely generated abelian group and let $S=\{s_1,\dots,s_m\}\subseteq G$ be a non-empty finite set.
Every point $a$ of the simplex
\[
A=\Bigl\{x\in\R^m:\ x_i\ge0,\ \textstyle\sum_i x_i=1\Bigr\}
\]
defines a homogeneous random walk on $G$, which moves from $g$ to $g+s_i$ with probability $a_i$. Li, Mossel and
Weiss~\cite[Section~6]{LMW} study the set $R\subseteq A$ of those $a$ for which this walk is recurrent, and its
complement $R^c=A\setminus R$ of transient walks. Identify $G$ with $\Z^n\times H$ for a finite abelian group $H$, let
$\pi:\Z^n\times H\to\Z^n\subseteq\R^n$ be the projection and put $v_i=\pi(s_i)$. Their Proposition~6.1 describes $R$:
\begin{equation}\label{eq:prop61}
a\in R\iff \dim\spn\{v_i:a_i>0\}\le2\quad\text{and}\quad \sum_{i=1}^m a_iv_i=0 .
\end{equation}
From this they prove that $R$ is closed~\cite[Theorem~6.2]{LMW}, that $R$ is path-connected when $S$ is
symmetric~\cite[Theorem~6.3]{LMW}, and that $R$ is convex if and only if all the spaces $\spn\{v_i:a_i>0\}$, $a\in R$,
lie in one subspace of dimension at most two~\cite[Theorem~6.5]{LMW}. Their Example~6.4 shows that $R$ can be
disconnected, and their Theorem~6.9 gives a sufficient condition for $R^c$ to be path-connected, which applies to every
symmetric $S$ whose image spans a space of dimension at least three~\cite[Corollary~6.10]{LMW}. They show that $R^c$
is not convex for symmetric $S$ unless every $v_i$ is zero~\cite[Proposition~6.12]{LMW}. They leave open the
classification of path-connectedness of $R$ for general $S$, and ask:
\begin{itemize}
\item[] Question~6.7. When is $R$ simply connected?
\item[] Question~6.8. When is $R^c$ path-connected?
\item[] Question~6.11. When is $R^c$ simply connected?
\item[] Question~6.13. When is $R^c$ convex?
\end{itemize}
They expected $R^c$ to be path-connected apart from one-dimensional examples such as $S=\{1,-1\}\subseteq\Z$, and the
answer to Question~6.11 to differ substantially from that of Question~6.8, on the strength of the example
$S=\{\pm e_1,\pm e_2\}\subseteq\Z^2$, where $R^c$ is path-connected but not simply connected.

The answers below depend only on the vectors $v_1,\dots,v_m$, which may repeat. We call a line $L$ through the origin
\emph{balanced} if it contains steps on both sides of $0$, that is, $v_i,v_j\in L$ with $v_j=-cv_i\ne0$ for some $c>0$.
We call a plane $V$ through the origin \emph{balanced} if it is spanned by the steps it contains and $0$ lies in the
convex hull of those steps. The \emph{balance graph} $B$ is the bipartite graph whose vertices are the balanced lines
and the balanced planes, a line being joined to each balanced plane that contains it. Let $I_0=\{i:v_i=0\}$ and
$r=\dim\spn\{v_1,\dots,v_m\}$. As in~\cite{LMW}, the empty set counts as path-connected and convex.

\begin{theorem}\label{thm:R}
\begin{enumerate}
\item If $I_0\ne\emptyset$, then $R$ is star-shaped, hence contractible.
\item If $r\le2$, then $R$ is convex.
\item If $I_0=\emptyset$ and $r\ge3$, then $R$ is homotopy equivalent to $B$.
\end{enumerate}
\end{theorem}

\begin{corollary}\label{cor:R}
$R$ is path-connected if and only if $I_0\ne\emptyset$, or $r\le2$, or $B$ is connected or empty. A non-empty $R$ is
simply connected if and only if $I_0\ne\emptyset$, or $r\le 2$, or $B$ is a tree. In case~(3) the components of $R$
correspond to the components of $B$, and the fundamental group of each component is free of rank equal to the number
of independent cycles of the corresponding component of $B$.
\end{corollary}

For symmetric $S$ with $r\ge3$ and $I_0=\emptyset$ the graph $B$ is connected, in line with~\cite[Theorem~6.3]{LMW},
but it always contains a cycle (Example~\ref{ex:sym}). So among symmetric $S$ the set $R$ is convex when $r\le2$ and not
convex when $r\ge3$~\cite[Corollary~6.6]{LMW}; when $r\ge3$ it is contractible if some $s_i$ projects to $0$, and path-connected but not simply connected
otherwise.

\begin{theorem}\label{thm:Rc}
Let $K=\conv\{v_1,\dots,v_m\}$.
\begin{enumerate}
\item If $r\ge3$, or $0\notin\relint K$, then $R^c$ is non-empty and star-shaped.
\item If $0\in\relint K$ and $r=1$, then $R^c$ has exactly two components, each of them star-shaped.
\item If $0\in\relint K$ and $r=2$, then $R^c$ is homotopy equivalent to a circle.
\end{enumerate}
The remaining case is $r=0$, in which every $v_i$ is zero, $R=A$ and $R^c$ is empty.
In particular $R^c$ is path-connected unless the $v_i$ span a line and lie on both sides of $0$, and a non-empty
$R^c$ is simply connected unless $0\in\relint K$ and $r\in\{1,2\}$.
\end{theorem}

Here $\relint K$ is the interior of $K$ relative to its affine hull; when $r\ge 1$ and $0\in \relint K$ the affine hull
of $K$ is $\spn\{v_i\}$. This confirms the expectation of~\cite{LMW} about Question~6.8 and makes precise the
difference from Question~6.11: the only extra obstruction to simple connectivity is the planar case, of which
$\{\pm e_1,\pm e_2\}$ is an instance. Part~(1) for $r\ge3$ contains~\cite[Theorem~6.9 and Corollary~6.10]{LMW}.

\begin{theorem}\label{thm:convex}
The following are equivalent.
\begin{enumerate}
\item $R^c$ is convex.
\item Every recurrent law is supported on $I_0$, that is, $R=\{a\in A: a_i=0\text{ for }i\notin I_0\}$.
\item No subspace of dimension one or two contains nonzero steps $v_i$ whose convex hull contains $0$.
\end{enumerate}
\end{theorem}

For symmetric $S$ condition~(3) fails as soon as some $v_i\ne0$, which recovers~\cite[Proposition~6.12]{LMW}. The set
in~(2) is the set $P$ of positive recurrent walks of~\cite[Remark~6.14]{LMW}, so~(2) says that $R^c$ is convex exactly
when every recurrent walk is positive recurrent.

The proofs are short once~\eqref{eq:prop61} is available. For Theorem~\ref{thm:R}(3) we cover $R$ by the polytopes on
which the walk lives in a fixed plane and apply the nerve theorem; the intersection pattern of these polytopes is
governed by the balanced lines, and the nerve is homotopy equivalent to $B$. For Theorem~\ref{thm:Rc} the case $r\ge3$ is a
star-shaped argument about an interior point of $A$, and the case $r\le2$ reduces, through a section of the drift map,
to the topology of $K\setminus\{0\}$.

\section{Notation}\label{sec:notation}

Throughout, $f:\R^m\to\R^n$ is the linear map $f(a)=\sum_ia_iv_i$, so $f(A)=K$. For a subspace $V\subseteq\R^n$ put
\[
C_V=\{a\in A:\ f(a)=0,\ a_i=0\text{ whenever }v_i\notin V\}.
\]
Each $C_V$ is a polytope (possibly empty), $C_V\subseteq C_W$ when $V\subseteq W$, and $C_V\cap C_W=C_{V\cap W}$
directly from the definition. By~\eqref{eq:prop61}, $a\in R$ if and only if $a\in C_V$ for some subspace $V$ of
dimension at most two; one may take $V=\spn\{v_i:a_i>0\}$. Only the set $\{i: v_i\in V\}$ matters, so finitely many
polytopes occur and $R$ is a finite union of polytopes. We write $F=\{a\in A: a_i=0\text{ for }i\notin I_0\}$ for the
face of $A$ spanned by the steps that project to $0$; then $C_{\{0\}}=F$.

\begin{lemma}\label{lem:nonempty}
Suppose $I_0=\emptyset$. A line $L$ has $C_L\ne\emptyset$ if and only if it is balanced. If $V$ is a plane spanned by
the steps it contains, then $C_V\ne\emptyset$ if and only if $V$ is balanced.
\end{lemma}

\begin{proof}
$C_V$ is non-empty exactly when $0$ is a convex combination of the steps lying in $V$. On a line with no zero steps
this means steps on both sides of $0$.
\end{proof}

\section{The recurrent set}\label{sec:R}

\begin{proof}[Proof of Theorem~\ref{thm:R}(1) and (2)]
(1) Let $z\in F$, which is non-empty, and $a\in R$. For $0<t<1$ the point $(1-t)a+tz$ has support
$\supp a\cup\supp z$, the steps added project to $0$, so $\spn\{v_i: ((1-t)a+tz)_i>0\}=\spn\{v_i:a_i>0\}$, and the
drift is $(1-t)f(a)+tf(z)=0$. By~\eqref{eq:prop61} the segment from $a$ to $z$ lies in $R$, so $R$ is star-shaped
about $z$.

(2) If $r\le2$, the dimension condition in~\eqref{eq:prop61} always holds, so $R=A\cap f^{-1}(0)$, an intersection of
convex sets.
\end{proof}

Part~(2) is the special case $V=\spn\{v_i\}$ of~\cite[Theorem~6.5]{LMW}, stated there as~\cite[Corollary~6.6]{LMW}.

For part~(3) assume from now on that $I_0=\emptyset$ and $r\ge3$. Let $\mathcal P$ be the set of balanced planes and
$\mathcal L$ the set of balanced lines.

\begin{lemma}\label{lem:cover}
$R=\bigcup_{V\in\mathcal P}C_V$. Every balanced line lies in at least one balanced plane.
\end{lemma}

\begin{proof}
Let $a\in R$ and $W=\spn\{v_i:a_i>0\}$, so $a\in C_W$ and $\dim W\le2$. Since $f(a)=0$ and no step is zero, $W\ne\{0\}$.
If $\dim W=1$, pick a step $v_j\notin W$, which exists because $r\ge3$, and enlarge $W$ to $W+\R v_j$. In either case
$a\in C_V$ for a plane $V$ spanned by steps, and $C_V\ne\emptyset$, so $V$ is balanced by Lemma~\ref{lem:nonempty}. The
same enlargement applied to a balanced line $L$ gives a plane spanned by steps that contains the steps of $L$, hence
$0$ in the convex hull of its steps.
\end{proof}

\begin{lemma}\label{lem:nerve}
Let $\mathcal Q\subseteq\mathcal P$ have at least two elements. Then $\bigcap_{V\in\mathcal Q}C_V\ne\emptyset$ if and
only if all planes of $\mathcal Q$ contain one common balanced line.
\end{lemma}

\begin{proof}
$\bigcap_{V\in\mathcal Q}C_V=C_U$ with $U=\bigcap_{V\in\mathcal Q}V$. Distinct planes meet in a line or in $\{0\}$,
so $U$ is a line or $\{0\}$. $C_{\{0\}}=F=\emptyset$ because $I_0=\emptyset$, and for a line $C_U\ne\emptyset$ exactly
when $U$ is balanced (Lemma~\ref{lem:nonempty}).
\end{proof}

Let $N$ be the nerve of the cover $\{C_V\}_{V\in\mathcal P}$: its vertices are the balanced planes, and a set of
planes spans a simplex when the corresponding polytopes have a common point. For $L\in\mathcal L$ let
$\mathcal P_L=\{V\in\mathcal P: L\subseteq V\}$, which is non-empty by Lemma~\ref{lem:cover}, and let $D_L$ be the full
simplex on $\mathcal P_L$. By Lemma~\ref{lem:nerve},
\begin{equation}\label{eq:N}
N=\bigcup_{L\in\mathcal L}D_L\ \cup\ \{\text{balanced planes containing no balanced line}\},
\end{equation}
and two distinct lines $L,L'$ lie together in at most one plane, so $D_L\cap D_{L'}$ is empty or a single vertex.

\begin{lemma}\label{lem:RN}
$R$ is homotopy equivalent to $N$.
\end{lemma}

\begin{proof}
Let $M=A\cap f^{-1}(0)$, a polytope. Each non-empty $C_V$ is obtained from $M$ by setting some coordinates to zero,
so it is a face of $M$, and so is every intersection of the $C_V$. Hence $R$ is a union of faces of $M$, and the
barycentric subdivision of the face complex of $M$ restricts to a triangulation of $R$ in which every $C_V$ is a
subcomplex. Every non-empty intersection of the $C_V$ is convex,
hence contractible, so by the nerve theorem for a CW complex covered by subcomplexes~\cite[Corollary~4G.3 and
Exercise~4 of Section~4.G]{Hat02}, $R\simeq N$.
\end{proof}

\begin{lemma}\label{lem:NB}
$N$ is homotopy equivalent to $B$.
\end{lemma}

\begin{proof}
For each balanced line $L$ add a new vertex $c_L$ and the cone $c_L*D_L$, the join of the point $c_L$ with $D_L$,
which is the full simplex on $\{c_L\}\cup\mathcal P_L$. Let $N''=N\cup\bigcup_L c_L*D_L$. The cones meet $N$ in the $D_L$, and two cones meet only
inside $N$.

Since $c_L*D_L$ and $D_L$ are both contractible, the inclusion $D_L\hookrightarrow c_L*D_L$ is a homotopy equivalence,
and a CW pair has the homotopy extension property~\cite[Proposition~0.16]{Hat02}, so $D_L$ is a deformation retract
of $c_L*D_L$~\cite[Corollary~0.20]{Hat02}. These retractions fix $N$, so together they deformation retract $N''$ onto
$N$.

Inside $c_L*D_L$ lies the star $T_L=c_L*\mathcal P_L$, the union of the edges from $c_L$ to the vertices of
$\mathcal P_L$. It is contractible, so by the same two results $c_L*D_L$ deformation retracts onto $T_L$, fixing
$T_L$ and in particular every vertex of $\mathcal P_L$. Two cones $c_L*D_L$ and $c_{L'}*D_{L'}$ meet in at most one
vertex of $\mathcal P$, which both retractions fix, so the retractions fit together and deformation retract $N''$
onto $\bigcup_L T_L$ together with the isolated vertices in~\eqref{eq:N}. That complex is $B$, with $c_L$ in the
role of the line $L$. Hence $N\simeq N''\simeq B$.
\end{proof}

Lemmas~\ref{lem:RN} and~\ref{lem:NB} prove Theorem~\ref{thm:R}(3). Corollary~\ref{cor:R} follows: a non-empty star-shaped
or convex set is contractible (the empty set is path-connected and convex by convention), a graph is path-connected exactly when it is connected, and a non-empty graph is simply
connected exactly when it is a tree; the fundamental group of a connected graph is free on its independent cycles.

\begin{example}[Li, Mossel and Weiss]\label{ex:64}
In $\Z^3$ take $S=\{e_3,\ e_1,\ -e_1+e_2,\ -e_1-e_2,\ -e_3\}$~\cite[Example~6.4]{LMW}. The only balanced line is
$\R e_3$; the balanced planes are the three planes spanned by $e_3$ and one of the other steps, and the plane
$x_3=0$, which contains no balanced line. So $B$ is a star with three edges together with an isolated vertex, and $R$
has two components, each contractible.
\end{example}

\begin{example}\label{ex:sym}
Let $S$ be symmetric with $I_0=\emptyset$ and $r\ge3$, and choose linearly independent steps $u_1,u_2,u_3$. The lines
$\R u_j$ are balanced, the planes $\R u_j+\R u_k$ are balanced and distinct, and
$\R u_1,\ \R u_1+\R u_2,\ \R u_2,\ \R u_2+\R u_3,\ \R u_3,\ \R u_3+\R u_1$ is a cycle in $B$. So $R$ is not simply
connected. For $S=\{\pm e_1,\dots,\pm e_n\}\subseteq\Z^n$, $n\ge3$, the graph $B$ is the complete graph on $n$ vertices
with every edge subdivided, so $R$ is homotopy equivalent to a wedge of $(n-1)(n-2)/2$ circles. For $n=3$ this is seen
directly: drift zero forces $a_{e_j}=a_{-e_j}$, and the dimension condition forces one of the three pairs to vanish,
so $R$ is the boundary of a triangle.
\end{example}

\section{The transient set}\label{sec:Rc}

\begin{proof}[Proof of Theorem~\ref{thm:Rc}]
\emph{The case $r\ge3$.} Here and below $\supp a=\{i:a_i>0\}$. Let $c$ be a point of the relative interior of $A$. For $x\in R^c$ and $0<t\le1$ the point
$(1-t)x+tc$ has every coordinate positive, so its steps span a space of dimension $r\ge3$, and it is transient
by~\eqref{eq:prop61}. Hence $R^c$ is star-shaped about $c$.

\emph{The case $r\le2$.} Now $R=A\cap f^{-1}(0)$, as in the proof of Theorem~\ref{thm:R}(2). Triangulate $K$ using
only points of $\{v_1,\dots,v_m\}$ as vertices, which is possible for the convex hull of any finite point set, and
for each vertex $w$ of the triangulation fix one index $i(w)$ with $v_{i(w)}=w$. On a
simplex with vertices $w_0,\dots,w_k$ write $x=\sum_j\lambda_jw_j$ in barycentric coordinates and set
$s(x)=\sum_j\lambda_je_{i(w_j)}$, where $e_i$ are the vertices of $A$. Since $i(w)$ depends only on the vertex $w$,
these definitions agree on common faces, so
$s:K\to A$ is continuous, and $f\circ s$ is the identity. The homotopy
\[
H(a,t)=(1-t)a+t\,s(f(a)),\qquad a\in A,\ 0\le t\le1,
\]
satisfies $f(H(a,t))=f(a)$, so it maps $R^c\times[0,1]$ into $R^c$, starts at the identity and ends in
$s(K\setminus\{0\})$, and it fixes $s(K\setminus\{0\})$ pointwise. Thus $R^c$ deformation retracts onto
$s(K\setminus\{0\})$, which is homeomorphic to $K\setminus\{0\}$ because $f\circ s$ is the identity.

If $0\notin K$, then $R=\emptyset$ and $R^c=A$ is convex. If $0\in K\setminus\relint K$, choose $p\in\relint K$ and
$c\in A$ with $f(c)=p$. For $a\in R^c$ and $0<t\le1$ the drift of $(1-t)a+tc$ is $(1-t)f(a)+tp$, which lies in
$\relint K$ and is therefore not $0$, so $(1-t)a+tc$ is transient. Hence $R^c$ is star-shaped about $c$.
If $0\in\relint K$ and $r\ge1$, then $K$ is a convex body of dimension $r$ in $\spn\{v_i\}$ with $0$ in its interior,
and radial projection from $0$ deformation retracts $K\setminus\{0\}$ onto the boundary of $K$, a sphere of dimension
$r-1$. For $r=1$ this is two points, so $R^c$ has two contractible components; for $r=2$ it is a circle. Finally, if
$r=0$ then $R=A$ and $R^c$ is empty.

In case~(2) each component is in fact star-shaped. If $v_i$ spans the line and $\sigma(i)\in\{+,-,0\}$ is the side of
$0$ on which $v_i$ lies, the two components are $\{a\in A: f(a)\text{ lies on the }+\text{ side}\}$ and the same with
$-$, and each is star-shaped about the barycentre of the face spanned by its own steps $\{i:\sigma(i)=\pm\}$.
\end{proof}

\begin{proof}[Proof of Theorem~\ref{thm:convex}]
(2)$\Rightarrow$(1). The set $F$ is a face of $A$. If $x,y\in A\setminus F$, some coordinate $x_i$ with $i\notin I_0$
is positive, so $((1-t)x+ty)_i>0$ for $0\le t<1$, and $y\notin F$; hence $A\setminus F$ is convex. Since $F\subseteq R$
always holds by~\eqref{eq:prop61}, condition~(2) says $R=F$.

(1)$\Rightarrow$(2). Suppose $x\in R$ has $x_i>0$ for some $i$ with $v_i\ne0$. Let $\Phi$ be the face of $A$ spanned by
$\supp x$, so $x$ lies in its relative interior. The linear map $f$ is not constant on $\Phi$, because the vertex $e_i$
of $\Phi$ has $f(e_i)=v_i\ne0=f(x)$. Choose a vector $d$ parallel to $\Phi$ with $f(d)\ne0$. For small $\varepsilon>0$
the points $x\pm\varepsilon d$ lie in $\Phi$ and have drift $\pm\varepsilon f(d)\ne0$, so both are transient, while
their midpoint $x$ is recurrent. Hence $R^c$ is not convex.

(2)$\Leftrightarrow$(3). If $x\in R$ charges a step $v_i\ne0$, let $J=\{j\in\supp x: v_j\ne0\}$. Since steps with
$v_j=0$ do not contribute to the drift, the normalized restriction of $x$ to $J$ is a convex combination of the nonzero
steps $v_j$, $j\in J$, equal to $0$, and these steps span a space $V$ of dimension one or two by~\eqref{eq:prop61}. So
(3) fails. Conversely, a convex combination of nonzero steps in such a $V$ that equals $0$ is a recurrent law
by~\eqref{eq:prop61} that charges a nonzero step, so (2) fails.
\end{proof}

\section{A computer check}\label{sec:checks}

Two programs in the repository~\cite{Repo} test the statements on random examples.

The first, \texttt{check\_nerve.py}, draws step sets with no zero step and $r\ge3$ and builds the cover $\{C_V\}$. It
decides every intersection of the polytopes by linear programming on their defining constraints (the HiGHS solver
in SciPy), not through
Lemma~\ref{lem:nerve}. It then computes the rational Betti numbers of the nerve from floating-point ranks of its boundary matrices,
whose entries are $0$ and $\pm1$, and compares them with the
number of components and of independent cycles of $B$. On 300 step sets in $\Z^3$ and $\Z^4$ it found no
disagreement; on Example~\ref{ex:64} and on $\{\pm e_j\}$ in $\Z^3$ and $\Z^4$ it gives the Betti numbers of two
points, a circle and a wedge of three circles. Random step sets rarely give a cycle in $B$, so
\texttt{check\_more.py} repeats the test on 200 step sets in $\Z^3$, $\Z^4$ and $\Z^5$ built from a few vectors
together with their negatives or doubled negatives. Of these, 138 have a cycle in $B$ and 181 have a balanced line in
three or more balanced planes, and again there was no disagreement.

The same program tests Theorem~\ref{thm:Rc} on 150 step sets in $\Z$, $\Z^2$ and $\Z^3$, with zero steps allowed and
every rank from zero to three. It lists the transient points of $A$ whose coordinates lie in $\frac1N\Z$, with $N=6$
or $8$. Two of them are joined when the open segment between them avoids $R$; on that segment the support is the union
of the two supports and the drift moves linearly, so the test reduces to whether $0$ lies on the open segment between
the two drifts, which is checked in exact rational arithmetic, and to the rank of the union of the supports, which is
computed in floating point. The number of components of the
resulting graph agreed with Theorem~\ref{thm:Rc} in all 150 cases. The runs above are
\texttt{python check\_nerve.py 300 1} and \texttt{python check\_more.py 200 150 7} (Python~3 with NumPy~2.0.2 and SciPy~1.13.1).

\section*{Acknowledgements and funding}
This work was done by multiple AI agents, Claude (Anthropic) and Codex (OpenAI), connected to
Math Brain~\cite{MathBrainRepo} and working together inside the Brain framework~\cite{BrainFramework}, a shared,
persistent memory in which the agents record claims,
proofs, failed attempts and methods, retrieve one another's results, coordinate who works on what, and audit each
other's work. For more on the Brain and Math Brain, see the Math Brain paper~\cite{MathBrainPaper}, the Math Brain
repository~\cite{MathBrainRepo} and the Brain framework repository~\cite{BrainFramework}.

The research presented in this paper was supported by the European Research Council (ERC) under the European Union's
Horizon Europe research and innovation programme (grant agreement No.~101041711), by the Simons Foundation as part of
the Collaboration on the Mathematical and Scientific Foundations of Deep Learning, by Heights Labs, by Convex Nexus
Capital, by the Israel Science Foundation (grant number 2258/19), and by the Israel Science Foundation (ISF Grant
4101/25). Computational resources were provided in part by Google Cloud through its education credits programme
(credit 530550103).

\section*{Declaration of competing interest}
The author holds equity in Heights Labs and in Convex Nexus Capital and has a management role at Convex Nexus Capital.
These entities are included in the funding declaration above.

\section*{Data availability}
The programs described in Section~\ref{sec:checks} are in the repository~\cite{Repo}.

\section*{Declaration of generative AI and AI-assisted technologies}
This paper is part of a broader research project on AI proving problems. Claude agents (Anthropic) proved
Theorems~\ref{thm:R}, \ref{thm:Rc} and~\ref{thm:convex}, wrote the programs of Section~\ref{sec:checks} and prepared
the manuscript; Codex agents (OpenAI) refereed the proofs. The setting, the characterization~\eqref{eq:prop61},
Example~\ref{ex:64} and the questions answered here are due to Li, Mossel and Weiss~\cite{LMW}. The new results are
the description of $R$ by the balance graph, the homotopy types of $R^c$ and the convexity criterion. The author
directed the work and is responsible for the paper.

\begin{thebibliography}{99}
\bibitem{Hat02} A.~Hatcher, \emph{Algebraic Topology}, Cambridge University Press, Cambridge, 2002.
\bibitem{LMW} R.~Li, E.~Mossel and B.~Weiss, \emph{Monotonicity, topology, and convexity of recurrence in random
walks}, Ann. Appl. Probab. \textbf{36} (2026), no.~5, 4291--4310, \doi{10.1214/26-AAP2329}. Numbered
statements are cited from \arxiv{2403.04108} (version 4, 2025).
\bibitem{Repo} Y.~Shmalo, \emph{The topology of the recurrent and transient sets: paper and check programs}, 2026,
\url{https://github.com/yspennstate/recurrent-set-topology}.
\bibitem{MathBrainRepo} Y.~Shmalo, \emph{Math Brain: public snapshot, code and tests}, 2026,
\url{https://github.com/yspennstate/math-brain-public/tree/8c9a6bd478c4ec4082347151b7ae16db50c917dd}.
\bibitem{MathBrainPaper} Y.~Shmalo, \emph{Math Brain: shared, auditable memory for multi-agent mathematical research},
2026, \doi{10.5281/zenodo.24150837},
\url{https://github.com/yspennstate/math-brain-public/tree/main/papers/math-brain-2026-10-24}.
\bibitem{BrainFramework} Y.~Shmalo, \emph{Brain Framework}, public software repository, 2026,
\url{https://github.com/yspennstate/brain-framework-public/tree/5629e644cdd510173ae97e9aa92642dabd22dc57}.
\end{thebibliography}

\end{document}
